\documentclass[a4paper, fleqn, 11pt]{article}

\input{CWT_Format}

% Additional packages for code listings
\usepackage{listings}
\usepackage{xcolor}

% Configure code listing style
\lstset{
    basicstyle=\ttfamily\small,
    breaklines=true,
    frame=single,
    framerule=0.8pt, % Adjust thickness of the border
    backgroundcolor=\color{smoke_white!10},
    numbers=left,
    numberstyle=\tiny\color{light_black},
    keywordstyle=\color{bubblegum_blue},
    commentstyle=\color{light_jade!75!light_black},
    stringstyle=\color{bubblegum_red},
    showstringspaces=false,
    rulecolor=\color{light_black} % Change this to the color you want for the box border
}

% Custom commands for technical terms
\newcommand{\pymodule}[1]{\texttt{#1}}
\newcommand{\pyfunction}[1]{\texttt{#1()}}
\newcommand{\filepath}[1]{\texttt{#1}}
\newcommand{\bashcmd}[1]{\texttt{#1}}

\newcommand{\sigone}[0]{***}
\newcommand{\sigfive}[0]{**\phantom{*}}
\newcommand{\sigten}[0]{*\phantom{**}}
\newcommand{\signada}[0]{\phantom{***}}
\newcommand{\sigse}[0]{\phantom{**}}

\begin{document}

\begin{center}
\emph{Online Supplement \red{(Not intended for publication)}} \\[2mm]
\Large{
\textbf{Nowcasting GDP Using Data Revisions in an Emerging Economy}} \\[5mm]

\normalsize
\iffalse
%\iftrue
\begin{tabular}{c c c  }
\textbf{ Jason Cruz} & \textbf{ Diego Winkelried}
&   \textbf{ Javier Torres} \\
\href{jj.cruza@up.edu.pe }{jj.cruza@up.edu.pe }  &\href{winkelried\_dm@up.edu.pe}{winkelried\_dm@up.edu.pe} &
\href{J.Torresgomez@up.edu.pe }{j.torresgomez@up.edu.pe }  \\[1mm]
\multicolumn{3}{c}{\emph{School of Economics and Finance, Universidad del Pac\'{i}fico (Lima, Per\'{u})}}
\end{tabular}\\[4mm]
\else
\vspace{1.5cm}
\fi
\today
\end{center}

This supplemental document provides detailed information on the data sources, collection procedures, processing steps, and computational implementation used to construct the first real-time dataset (RTD) of Peruvian GDP growth. It focuses on the methodological workflow, quality-control standards, and replication resources that underpin the empirical analysis.

In Peru, the National Institute of Statistics and Informatics (\href{https://www.inei.gob.pe/estadisticas/indice-tematico/economia/}{INEI}) produces and disseminates GDP statistics, while the Central Reserve Bank of Peru (\href{https://www.bcrp.gob.pe/publications/weekly-reports/weekly-report.html}{BCRP}) republishes these figures in its \href{https://www.bcrp.gob.pe/statistics/weekly-report-tables.html}{Weekly Report} (WR). Although INEI is the primary source, the BCRP provides a clearer and more consistent record of releases, reflecting legal obligations under article 84 of the Peruvian Constitution and articles 2 and 74 of the BCRP's Organic Law, which require the periodic publication of key national macroeconomic statistics. Our data collection strategy draws on both sources to construct the RTD of Peruvian GDP growth.

%-----------------------------------------------------------------------
% 1. Revision Calendar
%-----------------------------------------------------------------------

\section{Revision Calendar}
\label{section:calendar}

A formal publication calendar governs the initial release of GDP growth figures, which are published with a two-month lag relative to the reference period across all frequencies. While this delay aligns with international practice, it underscores the need to account for data revisions when using GDP for real-time policy analysis. In contrast, no official revision calendar is maintained. This distinction is important to highlight: the timing of first releases is regulated and predictable, whereas subsequent revisions follow no explicit or public schedule.

For the construction of a RTD, however, a revision calendar is indispensable. Without a consistent schedule of updates, it becomes impossible to define vintages and track revisions systematically. A regular calendar ensures comparability across time, while irregular release dates complicate measurement, since the horizon between releases varies.

To address this limitation, we reconstructed an implicit revision calendar using the BCRP's WR. Each WR reproduces INEI's figures and specifies the date of the update. We then harmonized this information to a monthly frequency—consistent with the ``horizon'' definition in the main paper—based on two criteria:

\begin{enumerate}
    \item \href{https://www.gob.pe/institucion/inei/normas-legales/2294897-316-2003-inei}{Chief Resolution No. 316-2003-INEI} mandates that sectoral offices update their data at least quarterly (e.g., March, June, September, December). These offices also produce monthly statistical information for the Monthly Index of National Production (IMPN), suggeting that revisions are made monthly.

    \item In the WR, revisions are effectively made on a month-to-month basis rather than week-to-week. Our reconstructed calendar confirms that the WR report updated GDP figures at least once per month.
\end{enumerate}

%-----------------------------------------------------------------------
% 2. Raw Data Sources
%-----------------------------------------------------------------------

\section{Raw Data Sources}
\label{section:raw_data}

Once the monthly revision calendar was established, we compiled the last WR of each month. Each monthly WR defines a \emph{vintage} of GDP growth rates, which together form the basis of the RTD.

The WR republishes GDP growth rates by sector as released by INEI. These sectoral growth rates are reported year-on-year and presented in a well-organized table whose structure has changed little over the past three decades, although the table numbering has varied across issues. Two key tables were systematically extracted: one reporting monthly growth rates, and another reporting quarterly and annual growth rates

Historically, until the early 2000s, the WR was titled the ``Weekly Bulletin'', and no digital versions of either the Bulletin or the Report exist prior to 2013. Consequently, for the period spanning 2000–2012, the macroeconomic tables are only accessible in hardcover volumes available at the BCRP's \href{https://biblioteca.bcrp.gob.pe/discovery/search?vid=51BCDRDP_INST:51BCDRDP2}{Renzo
 Rossini Library}, requiring in-person consultation. These volumes cannot be removed from the library, but they may be photocopied or photographed. This process is not always straightforward, as the historical WRs are often bound in large, heavy volumes.

From 2013 to the present, however, WRs have been freely accessible on the BCRP's website in PDF format, which has greatly facilitated data collection. The data are also presented in spreadsheets and in the BCRP data database, although only in their revised form.

\subsection{Scanned Data (pre–2013)}

For the pre-2013 period, the WRs  were scanned into PDF format, producing non-structured hardcopies. Consequently, the extraction process was far from trivial and followed three main steps: (i) Computer-based extraction, (ii) manual correction and verification, and (iii) construction of parsed structured tables feeding into automated vintage-building routines implemented in the \pymodule{peru\_gdp\_rtd.cleaners.old\_table\_cleaner} module.

The scanned sources posed significant challenges:
%
\begin{itemize}
    \item Warped or distorted text and misaligned table lines, often caused by hardcover bindings during scanning.
    \item Shadowing, opacity, or blur reducing legibility.
    \item Fragmented or incomplete tables.
    \item Highlighted entries interfering with optical recognition accuracy.
    \item Inconsistent font sizes across tables within the same report.
    \item Residual marks obscuring decimal points or digits.
    \item Ink and other physical spots disrupting recognition.
\end{itemize}

Figures~\ref{fig:wr_2_2007_table_1} and \ref{fig:wr_2_2007_table_2} show examples of scanned tables exhibiting some of these difficulties.

\begin{figure}[!t]
\caption{Table 70, WR 2 (January 2007): monthly growth rates}
\includegraphics[width=1\textwidth]{pictures/wr_2_2007_table_1.pdf}
\label{fig:wr_2_2007_table_1}\\
\footnotesize \textbf{Notes:} The low resolution, uneven lighting, and shadowing from the page fold in this figure illustrate common difficulties for OCR processing. Distortions in the horizontal and vertical alignment of cells, faint characters, and overlapping text reduce recognition accuracy. Highlighted values and residual marks (e.g., spots or smudges) can obscure digits or decimal points, leading to frequent misclassifications. In this specific case, blurred numbers in the central columns and shading near the binding make it particularly challenging to extract reliable machine-readable data without extensive manual correction.
\end{figure}



Given these limitations, data extraction was particularly demanding and did not provide definite results. We employed Optical Character Recognition (OCR) using the open-source \href{https://github.com/tesseract-ocr/tesseract}{\texttt{Tesseract}} engine in \texttt{Python}. Originally developed by Hewlett-Packard and now maintained by Google, \texttt{Tesseract} is widely used for digitizing printed material, including statistical reports.

\newpage
The OCR workflow consisted of:
%
\begin{enumerate}
    \item Converting pages into PNG format.
    \item Transforming visual text into machine-readable formats, with extensive pre-processing to improve accuracy:
        \begin{itemize}
            \item Grayscale conversion and binarization to separate text from background.
            \item Noise removal (e.g., speckles, smudges) using median filtering or adaptive thresholding.
            \item Skew correction (deskewing) to align rotated text.
            \item Image scaling to ensure at least 300 DPI resolution.
            \item Contrast enhancement (e.g., via CLAHE) to make characters more prominent.
            \item Thinning and skeletonization to separate thick or touching characters.
            \item Border removal to avoid false detections.
        \end{itemize}
    \item Extracting tabular text into editable formats.
\end{enumerate}

Although OCR provided an initial transcription, frequent misclassifications occurred (e.g., reading \texttt{9.4} as \texttt{5.4}), making manual verification indispensable.

\begin{figure}[!t]
\caption{Table 88, WR 2 (January 2007): quarterly and annual growth rates}  \label{fig:wr_2_2007_table_2}
\includegraphics[width=1\textwidth]{pictures/wr_2_2007_table_2.pdf}\\
\footnotesize \textbf{Notes:} Another example of a problematic scanned WR table. OCR processing is hindered by uneven contrast, faint print, shading near the binding, and inconsistent alignment. Highlighted values, small font sizes, and residual marks further obscure digits and decimal points, making reliable digitization difficult without extensive pre-processing and manual checks.
\end{figure}

Once raw tables were extracted and converted to readable CSV format, we processed the data using the \pymodule{peru\_gdp\_rtd.cleaners.old\_table\_cleaner} module and orchestrated the workflow through \pyfunction{old\_table\_1\_runner} and \pyfunction{old\_table\_2\_runner} functions. The pipeline involved:
%
\begin{enumerate}
    \item Data cleaning and validation using 70+ specialized transformation functions.
    \item Vintage-format conversion with target-event column structure.
    \item Concatenation across vintages to form complete RTD.
    \item Storage in both CSV and Parquet formats for accessibility.
\end{enumerate}

The combination of automated routines and manual oversight was essential for ensuring the reliability of the early vintages.



\subsection{Digital Data (post-2013)}

Since 2013, WRs have been published online as digital PDFs (semi-structured data), greatly facilitating automated extraction. We developed modular \texttt{Python} routines to handle heterogeneous table structures across reports. The primary extraction library was \texttt{tabula}, selected for its robustness in dealing with complex table layouts. All processing is implemented in the \pymodule{peru\_gdp\_rtd.cleaners.new\_table\_cleaner} module, with workflow orchestration through \pyfunction{new\_table\_1\_runner} and \pyfunction{new\_table\_2\_runner} functions.

\newpage
Despite their advantages, digital sources also presented several recurring issues:
%
\begin{itemize}
    \item Encoded characters replacing numerical values (e.g., \texttt{(cid:101)(cid:115)}).
    \item Incomplete vertical lines separating columns, creating ambiguous parsing.
    \item Duplicate documents (e.g., \texttt{ns-21-2015} and \texttt{ns-45-2015}), disrupting scraping routines.
    \item Empty table cells misinterpreted as headers (e.g., \texttt{ns-25-2014}, \texttt{ns-14-2020}).
    \item Missing target-period labels, complicating the assignment of growth rates (e.g., \texttt{ns-28-2013}).
    \item Missing year labels, complicating the chronological alignment of data (e.g., \texttt{ns-28-2013}).
    \item Duplicated or misaligned headers generating spurious columns.
    \item Double or fragmented border lines, complicating table parsing.
\end{itemize}

Figure~\ref{fig:wr_28_2013_table_1} shows an example of a monthly table in digital format, where the fifth issue above is apparent: missing year labels (from May to December). While it is obvious to a human reader that the missing label corresponds to year \texttt{2012}, automated scripts must explicitly account for such cases. On the other hand, Figure~\ref{fig:wr_48_2016_table_2} illustrates additional issues such as incomplete borders and duplicated headers.


\begin{figure}[t!]
\caption{Table 62, WR 28 (July 2013): monthly growth rates}\label{fig:wr_28_2013_table_1}

\includegraphics[width=1\textwidth]{pictures/wr_28_2013_table_1.pdf}\\
\footnotesize \textbf{Notes:} Although OCR accuracy is much higher than for scanned sources, gaps remain in the archive. For instance, the year 2012 is missing. Such discontinuities complicate the construction of a complete real-time dataset and require cross-checking against alternative sources.
\end{figure}

\begin{figure}[t!]
\caption{Table 82, WR 48 (December 2016): quarterly and annual growth rates}\label{fig:wr_48_2016_table_2}
\includegraphics[width=1\textwidth]{pictures/wr_48_2016_table_2.pdf}\\
\footnotesize \textbf{Notes:} While free of resolution problems, this table illustrates other challenges for automated extraction, such as irregular table borders, as the border to the right of ``I'' in year 2016. These features can disrupt parsing routines and require additional cleaning before integration into the RTD.
\end{figure}

\newpage
The replication codes were designed to systematically address these problems through modular, configuration-driven architecture. All cleaning operations avoid hardcoding, ensuring adaptability to format changes and preserving consistency across vintages. The modular design consists of 70+ specialized functions organized into seven categories: text normalization (4 functions), general table cleaning (22 operations), column handling (14 functions), and table-specific transformations (13 functions each for Tables 1 and 2).

The automated extraction workflow, implemented through the \bashcmd{python scripts/update\_rtd.py} command, executes the following steps:
%
\begin{enumerate}
    \item Automated downloading of WR PDFs from the BCRP website via Selenium-based web scraper (\pymodule{peru\_gdp\_rtd.scrapers.bcrp\_scraper}).
    \item Identification of relevant tables using keyword searches (\pymodule{peru\_gdp\_rtd.processors.pdf\_processor}).
    \item Extraction of tables with \texttt{tabula} and cleaning using specialized functions.
    \item Restructuring into standardized vintage format with type-safe configuration management.
    \item Concatenation across vintages using union-based column alignment.
    \item Export to multiple formats (CSV, Parquet) with comprehensive metadata.
\end{enumerate}

\newpage
Once both scanned and digital vintages were processed, they were integrated into a continuous RTD spanning 2000–2024. The integration workflow includes:
%
\begin{enumerate}
    \item Concatenating vintages from pre- and post-2013 samples (\pymodule{peru\_gdp\_rtd.transformers.concatenator}).
    \item Extracting revision metadata from PDF sources (\pymodule{peru\_gdp\_rtd.transformers.metadata\_handler}).
    \item Applying base-year mapping to identify structural breaks in GDP series.
    \item Generating sentinel-adjusted datasets to mark invalid comparisons due to base-year changes.
    \item Producing benchmark indicator datasets to flag methodological revisions.
    \item Converting vintage format to release format for econometric analysis (\pymodule{peru\_gdp\_rtd.transformers.releases\_converter}).
    \item Storage of all dataset variants with comprehensive documentation.
\end{enumerate}






%-----------------------------------------------------------------------
% 3. The Real-Time Dataset (RTD)
%-----------------------------------------------------------------------

\section{The Real-Time Dataset (RTD)}
\label{section:rtd}
We emphasize the distinction between two key products:

\begin{itemize}
    \item \textbf{Vintages dataset:} contains all successive records as originally extracted. This constitutes the real-time dataset (RTD).
    \item \textbf{Releases dataset:} derived from vintages, reorganized for revision analysis, and serving as the main input for our empirical tests.
\end{itemize}

The vintages dataset is the direct output of the extraction process and preserves the chronology of publications exactly as they appeared. By contrast, the releases dataset is a research-oriented transformation of the RTD, specifically designed to facilitate the study of revisions. Both formats are complementary: the vintages dataset ensures transparency and replicability, while the releases dataset provides the structure needed for econometric testing.

To illustrate the distinction between vintages and releases, we focus on a window covering the COVID-19 lockdown period, during which monthly GDP growth rates exhibit large negative values that are visually salient. Table~\ref{tab:vintages} presents the real-time dataset (RTD) in vintage format, where columns correspond to publication vintages (e.g., 2020m05, 2020m06, \ldots, 2021m05) and rows correspond to target periods (e.g., 2020m03, 2020m04, etc.). Each cell reports the value for a given target period as published in a specific vintage. To facilitate visual tracking across vintages, first releases are highlighted in blue, second releases in red, third releases in green, and fourth releases in yellow within each target-period row. For example, in the 2020m07 vintage column, the third estimate of March 2020 growth is collected (target period 2020m03 highlighted in green, $-16.30$), while the first estimate for May 2020 is collected (target period 2020m05 highlighted in blue, $-32.80$).

The releases dataset reorganizes the same underlying information by changing the unit of organization. Table~\ref{tab:releases} presents this alternative structure, where rows still correspond to target periods, but columns now group observations by release sequence rather than by vintage. Specifically, the column labeled First release ($y_t^1$) collects the first published estimate for each target period, Second release ($y_t^2$) collects the second estimate, and so on, up to the Last release ($y_t^H$). This mapping can be seen directly by tracking highlighted cells across tables. For instance, the blue cells corresponding to the first estimates of 2020m03 ($-16.30$), 2020m05 ($-32.80$), and 2020m07 ($-11.70$) appear in different vintage columns in Table~\ref{tab:vintages}, but they all appear together in the $y_t^1$ column of Table~\ref{tab:releases}. Likewise, the red cells (second releases) for the same target periods are grouped under $y_t^2$, the green cells under $y_t^3$, and the yellow cells under $y_t^4$.

This reorganization---standard in the real-time data literature---makes the revision process explicit and operational. Revisions are computed as differences between consecutive release columns (e.g., $y_t^2 - y_t^1$, $y_t^3 - y_t^2$). The Last release column ($y_t^H$) contains the most recent available estimate for each target period, after which no further economically meaningful revisions are observed. By making the flow of information updates transparent, the releases format provides a clean and interpretable structure for the rationality, efficiency, and benchmarking analyses conducted in the main paper.

\begin{table}[t!]
\caption{Dataset of monthly vintages}\label{tab:vintages}
\resizebox{\textwidth}{!}{%
\input{tables/Tab_Vintages}
}\\[2mm]
\footnotesize \textbf{Notes:} This table illustrates the vintage format of a real-time database. Each column corresponds to a specific data vintage — that is, the set of estimates available in the WR at the end of a given month. Each row corresponds to a target period (e.g., a specific month for which GDP growth is estimated). The publication lag of the first release is two months.
\end{table}


\begin{table}[!b]
\caption{Dataset of monthly releases}
\label{tab:releases}
\resizebox{\textwidth}{!}{\footnotesize{%
\input{tables/Tab_Releases}
}
%
}\\[2mm]
\footnotesize \textbf{Notes:} This table presents the release format derived directly from the vintage representation in Table~\ref{tab:vintages}. The data are not modified; they are only reorganized to group estimates by release order rather than by publication vintage. Each row corresponds to a target period (e.g., a specific month of GDP growth), and each column reports $y_t^h$, the $h$-th published release of that target-period estimate.
\end{table}

%-----------------------------------------------------------------------
% 4. Coverage of the RTD
%-----------------------------------------------------------------------
\clearpage
\section{Coverage of the RTD}
\label{section:gdp_rtd_indicators}

We now summarize the coverage and structure of the RTD, which supports a comprehensive analysis of GDP as a real-time indicator of economic activity.

The dataset comprises monthly, quarterly, and annual GDP growth series for \emph{aggregate real GDP} and for \emph{all industries covered by the project}, with temporal coverage and revision horizons reported in Table~\ref{tab:rt_sum}. Specifically, the RTD includes the following industries: agriculture and livestock, fishing, mining and fuel, manufacturing, electricity and water, construction, commerce, and other services, in addition to total real GDP.

All industry-level series correspond to INEI’s level-9 classification, the most disaggregated level available in the official national accounts. Table~\ref{tab:rt_sum} reports the complete and exhaustive set of industries included in the RTD, together with their corresponding frequencies, target-period coverage, and revision horizons. No additional industry series beyond those listed in the table are included in the project.

Table~\ref{tab:rt_sum} reports the complete and exhaustive set of sectors included in the RTD, together with their corresponding frequencies, target-period coverage, and revision horizons. No additional sectoral series beyond those listed in the table are included in the project.

\begin{table}[b!]
\caption{Sectoral GDP as a real-time indicator}
\label{tab:rt_sum}
\resizebox{\columnwidth}{!}{
\input{tables/Tab_RT_Summarize}
%
}\\[2mm]
\footnotesize \textbf{Notes:} The disaggregation of economic activities follows level 9 of INEI's official classification. Under this classification, the ``other services'' sector includes a diverse set of activities, such as real estate and personal services, including education and health (14.89\%). The following categories are also identified: transportation, storage, postal and courier services (4.97\%); accommodation and food services (2.86\%); telecommunications and other information services (2.66\%); financial and insurance services (3.22\%); business services (4.24\%); product taxes (8.29\%); and public administration and defense (4.29\%).
\end{table}

\clearpage

%-----------------------------------------------------------------------
% 5. Replication and Availability
%-----------------------------------------------------------------------

\section{Computational Implementation and Replication}
\label{section:replication}

We provide a production-ready, open-source pipeline for constructing and updating the Peruvian GDP Real-Time Dataset. The codebase comprises over 4,000 lines of professionally documented Python code, organized into 14 specialized modules following software engineering best practices. All code is publicly available on GitHub (\url{https://github.com/JasonCruz18/peruvian_gdp_revisions}) under an MIT license, ensuring full transparency and reproducibility.

The pipeline architecture ensures complete reproducibility through configuration-driven execution, idempotent operations, and comprehensive error handling. A single command (\bashcmd{python scripts/update\_rtd.py}) executes the entire 6-step workflow from raw PDF download to final dataset generation, producing 10+ dataset variants suitable for econometric analysis. The modular design facilitates future extensions and adaptations to changes in data source formats while maintaining code maintainability and testability.

\subsection{Software Architecture}

The codebase follows a modular architecture with clear separation of concerns, organized into seven primary modules:

\paragraph{Configuration Module (\pymodule{peru\_gdp\_rtd.config}):} Manages all pipeline settings through a 231-line YAML configuration file (\filepath{config/config.yaml}) and type-safe Python dataclasses. This design eliminates hardcoding entirely, enabling users to customize paths, scraping parameters, cleaning tolerances, and output formats without modifying source code. The configuration system uses Pydantic-style validation to ensure type safety and catch errors early.

\paragraph{Scrapers Module (\pymodule{peru\_gdp\_rtd.scrapers}):} Implements automated web scraping from the BCRP website using Selenium WebDriver. The \pyfunction{pdf\_downloader} function handles authentication, rate limiting, retry logic with exponential backoff, and concurrent downloads while respecting server policies. The scraper automatically organizes PDFs by year and maintains download records for idempotent execution.

\paragraph{Processors Module (\pymodule{peru\_gdp\_rtd.processors}):} Handles PDF processing tasks including keyword-based page extraction (\pyfunction{pdf\_input\_generator}), file organization (\pymodule{file\_organizer}), and metadata parsing (\pymodule{metadata}). These functions reduce raw 4-page Weekly Reports to 2-page inputs containing only GDP tables, significantly improving downstream processing efficiency.

\paragraph{Cleaners Module (\pymodule{peru\_gdp\_rtd.cleaners}):} Contains 70+ specialized data cleaning functions organized into seven sub-modules:
\begin{itemize}
    \item \pymodule{text\_cleaners} (4 functions): Text normalization, whitespace handling, special character removal
    \item \pymodule{table\_cleaners} (22 functions): General DataFrame operations, noise removal, validation
    \item \pymodule{column\_handlers} (14 functions): Column renaming, type conversion, standardization
    \item \pymodule{table1\_cleaners} (13 functions): Monthly table-specific transformations
    \item \pymodule{table2\_cleaners} (13 functions): Quarterly/annual table-specific transformations
    \item \pymodule{old\_table\_cleaner} (1 class): Orchestrates cleaning for scanned CSV data (pre-2013)
    \item \pymodule{new\_table\_cleaner} (1 class): Orchestrates cleaning for digital PDF data (post-2013)
\end{itemize}

Each function is documented with type hints, docstrings, and example usage, ensuring transparency in data transformations.

\paragraph{Transformers Module (\pymodule{peru\_gdp\_rtd.transformers}):} Implements data transformation and RTD construction through four specialized modules:
\begin{itemize}
    \item \pymodule{vintage\_preparator}: Reshapes cleaned tables into vintage format with target-period columns
    \item \pymodule{concatenator}: Merges vintages across years using union-based column alignment (372 lines)
    \item \pymodule{metadata\_handler}: Extracts revision metadata, applies base-year mappings, generates benchmark indicators (655 lines)
    \item \pymodule{releases\_converter}: Transforms vintage format to release format using NumPy-based vertical alignment (241 lines)
\end{itemize}

\paragraph{Orchestration Module (\pymodule{peru\_gdp\_rtd.orchestration}):} Provides high-level workflow coordination through four runner functions (\pyfunction{old\_table\_1\_runner}, \pyfunction{old\_table\_2\_runner}, \pyfunction{new\_table\_1\_runner}, \pyfunction{new\_table\_2\_runner}) that orchestrate complete processing pipelines from raw data to structured vintages.

\paragraph{Utilities Module (\pymodule{peru\_gdp\_rtd.utils}):} Contains shared infrastructure including \pymodule{RecordManager} for idempotent operations and \pymodule{alerts} for user notifications.

\subsection{Installation and Setup}

The pipeline requires Python 3.10 or higher and Java Runtime Environment (JRE) for PDF processing. Installation follows standard Python packaging conventions:

\begin{lstlisting}[language=bash, caption={Installation procedure}]
# Clone repository
git clone https://github.com/JasonCruz18/peruvian_gdp_revisions.git
cd peru_gdp_revisions

# Install Python dependencies (pip)
pip install -r requirements.txt

# Alternative: Install with Conda
conda env create -f environment.yml
conda activate gdp_revisions

# Verify installation
python -c "from peru_gdp_rtd.config import get_settings; print('OK')"
\end{lstlisting}

Dependencies include:
\begin{itemize}
    \item \texttt{pandas} (2.0+): DataFrame operations
    \item \texttt{numpy} (1.24+): Numerical computations
    \item \texttt{selenium} (4.0+): Web scraping
    \item \texttt{tabula-py} (2.5+): PDF table extraction
    \item \texttt{PyMuPDF} (1.23+): PDF metadata extraction
    \item \texttt{tqdm} (4.65+): Progress tracking
    \item \texttt{pyyaml} (6.0+): Configuration parsing
\end{itemize}

The \filepath{requirements.txt} file pins exact versions for reproducibility.

\subsection{Pipeline Execution}

The pipeline supports both complete workflow execution and selective step execution through a command-line interface:

\begin{lstlisting}[language=bash, caption={Pipeline execution examples}]
# Execute complete pipeline (Steps 1-6)
python scripts/update_rtd.py

# Execute specific steps only
python scripts/update_rtd.py --steps 3,4,5,6  # Skip download/extraction

# Skip PDF download (useful for testing)
python scripts/update_rtd.py --skip-download

# Verbose output for debugging
python scripts/update_rtd.py --verbose

# Dry run (show execution plan without running)
python scripts/update_rtd.py --dry-run

# Use custom configuration file
python scripts/update_rtd.py --config path/to/custom_config.yaml
\end{lstlisting}

The six pipeline steps are:

\begin{enumerate}
    \item \textbf{Download PDFs}: Automated scraping from BCRP website (2013--present)
    \item \textbf{Generate Inputs}: Extract GDP table pages from full PDFs
    \item \textbf{Clean \& Build RTD}: Process both OLD (CSV) and NEW (PDF) sources into vintage format
    \item \textbf{Concatenate}: Merge vintages across all years into unified RTD
    \item \textbf{Metadata}: Extract revision information, apply base-year mappings, generate benchmarks
    \item \textbf{Convert to Releases}: Transform vintage format to release format for analysis
\end{enumerate}

Each step maintains processing records to enable idempotent execution—rerunning the pipeline only processes new or modified data, not previously completed work.

\subsection{Data Cleaning Pipeline}

The data cleaning workflow applies 70+ specialized transformations to handle the formatting irregularities described in Section~\ref{section:raw_data}. Key cleaning operations include:

\paragraph{Text Normalization:}
\begin{itemize}
    \item Unicode normalization (NFC form)
    \item Whitespace standardization (tabs → spaces, multiple spaces → single)
    \item Special character handling (accented characters, punctuation)
    \item Case normalization for sector names
\end{itemize}

\paragraph{Table Structure Cleaning:}
\begin{itemize}
    \item Empty row removal (fully NA or whitespace-only)
    \item Duplicate row detection and removal
    \item Header row identification and extraction
    \item Column alignment correction (handling irregular borders)
    \item Numeric column identification (distinguishing from text columns)
\end{itemize}

\paragraph{Numerical Data Cleaning:}
\begin{itemize}
    \item Encoded character replacement (e.g., \texttt{(cid:101)} → valid numbers)
    \item Decimal separator normalization (comma → period)
    \item Thousands separator removal
    \item Type conversion with error handling (string → float)
    \item Missing value handling (empty strings, dashes, ``n.d.'' → NA)
\end{itemize}

\paragraph{Temporal Data Handling:}
\begin{itemize}
    \item Missing year label inference (forward-filling from previous columns)
    \item Month abbreviation parsing (Spanish: ``ene'', ``feb'', etc.)
    \item Quarter identification (I, II, III, IV → 1, 2, 3, 4)
    \item Target period construction (year + month/quarter → ``2020m5'', ``2020q2'')
\end{itemize}

\paragraph{Sector Name Standardization:}
The pipeline maintains comprehensive mappings for sector name variations, handling both Spanish and English terminology:

\begin{lstlisting}[language=Python, caption={Example: Sector name standardization}, label={lst:sector}]
SECTOR_MAP = {
    "agriculture and livestock": "agriculture",
    "agropecuario": "agropecuario",
    "mining and fuel": "mining",
    "mineria e hidrocarburos": "mineria e hidrocarburos",
    "manufacturing": "manufacturing",
    "manufactura": "manufactura",
    # ... 15+ additional mappings
}
\end{lstlisting}

These mappings ensure consistency across vintages despite variations in source formatting.

\subsection{Configuration Management}

All pipeline behavior is controlled through \filepath{config/config.yaml}, a human-readable configuration file with no hardcoded values in source code. Key configuration categories include:

\begin{lstlisting}[language=yaml, caption={Configuration file excerpt (config/config.yaml)}]
# Scraper configuration
scraper:
  browser: "chrome"
  headless: false
  max_downloads: 60
  min_wait: 5.0
  max_wait: 10.0
  css_selectors:
    report_list: "#rightside ul.listado-bot-std-claros"

# PDF processing configuration
pdf_processing:
  keywords:
    - "ECONOMIC SECTORS"
    - "SECTORES ECONOMICOS"

# Data cleaning configuration
cleaning:
  decimal_places: 1
  pipeline_version: "1.0.0"
  sector_mappings_english:
    "agriculture and livestock": "agriculture"
    "fishing": "fishing"
    # ...

# Metadata configuration
metadata:
  base_years:
    - year: 2000
      wr: 28
      base_year: 1994
    - year: 2014
      wr: 11
      base_year: 2007
    - year: 2022
      wr: 20
      base_year: 2019

# Output configuration
features:
  persist_format: "csv"  # or "parquet"
  validate_data: true
\end{lstlisting}

This design allows researchers to adapt the pipeline to new data sources or methodological requirements without modifying Python code.

\subsection{Quality Assurance Features}

The pipeline implements several quality assurance mechanisms to ensure reliability and reproducibility:

\paragraph{Idempotent Operations:} The \pymodule{RecordManager} class tracks all processed files using text-based record files. Rerunning the pipeline automatically skips previously processed data, enabling incremental updates. Record files are human-readable and version-controlled for transparency.

\paragraph{Progress Tracking:} All long-running operations display \texttt{tqdm} progress bars with estimated completion times, colored output for status indicators, and summary statistics upon completion. This provides real-time feedback during execution and helps identify performance bottlenecks.

\paragraph{Error Resilience:} The pipeline implements graceful error handling with three-tier recovery:
\begin{enumerate}
    \item \textbf{Skip \& Continue}: Individual file failures don't halt entire workflow
    \item \textbf{Error Logging}: Failed operations are logged with full context for debugging
    \item \textbf{Summary Reports}: End-of-step summaries report success/failure counts
\end{enumerate}

\paragraph{Data Validation:} Built-in validation checks include:
\begin{itemize}
    \item Schema validation (expected columns present, correct types)
    \item Range validation (growth rates within plausible bounds)
    \item Consistency checks (monotonicity of vintage dates)
    \item Completeness checks (missing value patterns)
\end{itemize}

\paragraph{Dual Persistence:} The pipeline supports both CSV (human-readable, universally compatible) and Parquet (compressed, type-preserving) formats. Output datasets include both formats by default, with CSV for inspection and Parquet for analysis.

\subsection{Output Datasets}

The pipeline generates 10+ dataset variants optimized for different analytical purposes:

\begin{table}[b!]
\caption{Output dataset specifications}
\label{tab:output_dataset}
\resizebox{\columnwidth}{!}{
\input{tables/Tab_Output_Data}
%
}\\[2mm]
%\footnotesize \textbf{Notes:}
\end{table}

\paragraph{Vintage Format Schema:}
\begin{itemize}
    \item \texttt{industry} (string): Sector name (standardized)
    \item \texttt{vintage} (string): Publication month in ``YYYYmM'' format
    \item \texttt{tp\_YYYYmM} or \texttt{tp\_YYYYqN} or \texttt{tp\_YYYY} (float): Growth rate columns for each target period
\end{itemize}

\paragraph{Release Format Schema:}
\begin{itemize}
    \item \texttt{target\_period} (string): Reference period in ``YYYYmM''/``YYYYqN''/``YYYY'' format
    \item \texttt{industry\_h} (float): $h$-th release for each industry (e.g., \texttt{agriculture\_1}, \texttt{agriculture\_2}, etc.)
\end{itemize}

\paragraph{Benchmark Datasets:}
Replace growth rates with binary indicators where 1.0 denotes a benchmark revision (simultaneous update of monthly and quarterly/annual tables) and 0.0 otherwise. These datasets enable isolation of methodological revisions from routine data updates.

\paragraph{Base-Year Adjusted Datasets:}
Apply sentinel value (-999999.0) to growth rates that become incomparable due to base-year changes. This marks invalid observations without removing them, preserving the complete revision history while flagging structural breaks.

\subsection{Computational Requirements}

\paragraph{Estimated Runtime:}
\begin{itemize}
    \item Complete pipeline (Steps 1--6): Approximately 45--60 minutes on standard hardware
    \item Steps 3--6 only (processing existing PDFs): Approximately 25--35 minutes
    \item Incremental updates (new vintages only): Approximately 5--10 minutes
\end{itemize}

Actual runtime depends on internet connection speed (Step 1), CPU performance (Steps 3--5), and disk I/O (Step 6).

\paragraph{Hardware Requirements:}
\begin{itemize}
    \item CPU: Multi-core processor recommended (parallel processing support planned)
    \item RAM: 4 GB minimum, 8 GB recommended for large datasets
    \item Disk: 2 GB for code and dependencies, 5 GB for full dataset archive
    \item Internet: Stable connection required for Step 1 (PDF download)
\end{itemize}

\paragraph{Software Requirements:}
\begin{itemize}
    \item Operating System: Windows, macOS, or Linux (tested on all three)
    \item Python: 3.10+ (3.11 recommended)
    \item Java: JRE 8+ for \texttt{tabula-py}
    \item Browser: Chrome, Firefox, or Edge for Selenium (Chrome recommended)
\end{itemize}

\subsection{Repository Structure}

The complete repository structure follows modern Python packaging conventions:

\begin{lstlisting}[caption={Repository directory structure}]
peru_gdp_revisions/
|-- peru_gdp_rtd/                 # Main Python package
|   |-- config/                   # Configuration management
|   |   |-- config.yaml           # YAML configuration (231 lines)
|   |   `-- settings.py           # Type-safe Settings classes
|   |-- scrapers/                 # Web scraping
|   |   `-- bcrp_scraper.py       # Selenium-based PDF downloader
|   |-- processors/               # PDF processing
|   |   |-- pdf_processor.py      # Keyword-based extraction
|   |   |-- file_organizer.py     # File organization
|   |   `-- metadata.py           # Metadata parsing
|   |-- cleaners/                 # Data cleaning (70+ functions)
|   |   |-- text_cleaners.py
|   |   |-- table_cleaners.py
|   |   |-- column_handlers.py
|   |   |-- table1_cleaners.py
|   |   |-- table2_cleaners.py
|   |   |-- old_table_cleaner.py
|   |   `-- new_table_cleaner.py
|   |-- transformers/             # Data transformation
|   |   |-- vintage_preparator.py
|   |   |-- concatenator.py       # RTD merging (372 lines)
|   |   |-- metadata_handler.py   # Metadata extraction (655 lines)
|   |   `-- releases_converter.py # Format conversion (241 lines)
|   |-- orchestration/            # Workflow coordination
|   |   `-- runners.py            # High-level orchestrators
|   `-- utils/                    # Shared utilities
|       |-- data_manager.py       # RecordManager class
|       `-- alerts.py             # User notifications
|-- scripts/
|   |-- update_rtd.py             # One-button execution script
|   `-- validate_rtd.py           # Data quality validation
|-- notebooks/                    # Tutorial Jupyter notebooks
|   |-- 01_web_scraping.ipynb     # BCRP scraping tutorial
|   |-- 02_pdf_processing.ipynb   # PDF extraction tutorial
|   |-- 03_data_cleaning.ipynb    # Cleaning pipeline tutorial
|   |-- 04_rtd_construction.ipynb # RTD building tutorial
|   |-- 05_metadata_management.ipynb # Metadata handling
|   |-- 06_releases_datasets.ipynb    # Release conversion
|   |-- new_gdp_rtd.ipynb         # Newer vintage processing
|   |-- old_gdp_rtd.ipynb         # Older vintage processing
|   `-- README.md                 # Notebooks documentation
|-- dashboard/                    # Interactive Streamlit dashboard
|   |-- app.py                    # Main dashboard application
|   |-- config.py                 # Dashboard configuration
|   |-- requirements.txt          # Dashboard dependencies
|   |-- assets/                   # Dashboard assets
|   `-- README.md                 # Dashboard documentation
|-- docs/                         # Comprehensive documentation
|   |-- INSTALLATION.md           # Installation guide
|   |-- USAGE.md                  # Usage guide with examples
|   |-- ARCHITECTURE.md           # Software architecture
|   |-- CONTRIBUTING.md           # Contribution guidelines
|   |-- DATA_AVAILABILITY.md      # Data availability statement
|   `-- DOI_ASSIGNMENT.md         # Zenodo DOI guide
|-- .github/                      # GitHub configuration
|   `-- workflows/                # CI/CD pipelines
|       |-- tests.yml             # Comprehensive test suite
|       `-- quick-test.yml        # Fast smoke tests
|-- config/
|   |-- config.yaml               # User configuration
|   `-- config.example.yaml       # Configuration template
|-- data/                         # Generated datasets (gitignored)
|   |-- input/                    # Intermediate data
|   |-- output/                   # Final RTD datasets
|   `-- records/                  # Processing records
|-- metadata/
|   `-- wr_metadata.csv           # Revision metadata (tracked)
|-- tests/                        # Test suite
|   `-- test_smoke.py             # Smoke tests for all modules
|-- assets/                       # Project assets
|   `-- README.md                 # Asset documentation
|-- CITATION.cff                  # Citation metadata (CFF format)
|-- FAQ.md                        # Frequently Asked Questions
|-- LICENSE                       # MIT License
|-- environment.yml               # Conda environment specification
|-- pyproject.toml                # Modern Python packaging
|-- requirements.txt              # Dependency specifications
|-- requirements-frozen.txt       # Pinned versions for reproducibility
`-- README.md                     # Project documentation
\end{lstlisting}

\subsection{Interactive Dashboard and Visualization}

To facilitate exploration of the RTD, we developed an interactive web-based dashboard using Streamlit. The dashboard provides real-time data visualization, filtering, and export capabilities without requiring programming knowledge. Key features include:

\paragraph{Data Exploration:}
\begin{itemize}
    \item Interactive selection of datasets (monthly, quarterly, annual)
    \item Sector-level filtering with multi-select capabilities
    \item Vintage-based temporal navigation
    \item Search functionality across all available data
\end{itemize}

\paragraph{Visualizations:}
\begin{itemize}
    \item \textbf{Heatmaps}: Real-time dataset visualization showing revision patterns across vintages
    \item \textbf{Time series}: Growth rate evolution with vintage overlays
    \item \textbf{Revision analysis}: Histogram and box-plot distributions of revisions
    \item \textbf{Summary statistics}: Descriptive tables and metrics by sector
\end{itemize}

\paragraph{Export Capabilities:}
\begin{itemize}
    \item CSV export for spreadsheet analysis
    \item Excel export with multiple sheets
    \item JSON export for programmatic access
    \item Filtered dataset downloads based on user selections
\end{itemize}

The dashboard is launched via \bashcmd{streamlit run dashboard/app.py} and automatically opens in the default web browser. All dashboard configuration is managed through \filepath{dashboard/config.py}, enabling customization without modifying the application code.

\subsection{Tutorial Notebooks and Documentation}

The repository includes six comprehensive Jupyter notebooks that provide step-by-step tutorials for each stage of the RTD construction pipeline:

\begin{enumerate}
    \item \filepath{01\_web\_scraping.ipynb}: Demonstrates automated downloading of Weekly Reports from the BCRP website using Selenium, including authentication handling and rate limiting.

    \item \filepath{02\_pdf\_processing.ipynb}: Shows PDF table extraction using \texttt{tabula-py}, keyword-based page filtering, and metadata parsing.

    \item \filepath{03\_data\_cleaning.ipynb}: Illustrates the 70+ cleaning functions with before/after examples, handling of encoded characters, and validation checks.

    \item \filepath{04\_rtd\_construction.ipynb}: Guides through vintage preparation, concatenation logic, and RTD assembly from individual vintages.

    \item \filepath{05\_metadata\_management.ipynb}: Explains base-year mapping, benchmark indicator generation, and sentinel value application for structural breaks.

    \item \filepath{06\_releases\_datasets.ipynb}: Demonstrates conversion from vintage to release format, including the NumPy-based vertical alignment algorithm.
\end{enumerate}

Each notebook contains:
\begin{itemize}
    \item Executable code cells with complete examples
    \item Markdown explanations of methodology
    \item Visual outputs (tables, charts) showing intermediate results
    \item Links to relevant code modules and functions
    \item Tips for troubleshooting common issues
\end{itemize}

Additionally, the \filepath{docs/} directory provides comprehensive markdown documentation:
\begin{itemize}
    \item \filepath{INSTALLATION.md}: Detailed installation instructions for Windows, macOS, and Linux
    \item \filepath{USAGE.md}: Complete usage guide with examples and command-line reference
    \item \filepath{ARCHITECTURE.md}: Software architecture documentation with module descriptions
    \item \filepath{CONTRIBUTING.md}: Guidelines for contributing to the project
    \item \filepath{DATA\_AVAILABILITY.md}: Data availability statement for journal submission
    \item \filepath{DOI\_ASSIGNMENT.md}: Step-by-step guide for obtaining Zenodo DOI
\end{itemize}

\subsection{Continuous Integration and Testing}

The repository implements a comprehensive CI/CD pipeline using GitHub Actions, ensuring code quality and cross-platform compatibility:

\paragraph{Test Suite:}
Automated smoke tests (\filepath{tests/test\_smoke.py}) verify:
\begin{itemize}
    \item Package importability across all modules
    \item Configuration file validity and type safety
    \item Path resolution across different repository names
    \item Sector and month mapping completeness
    \item Base-year configuration consistency
\end{itemize}

\paragraph{Cross-Platform Testing:}
The CI pipeline (\ filepath{.github/workflows/tests.yml}) executes tests on:
\begin{itemize}
    \item \textbf{Operating Systems}: Ubuntu (latest), Windows (latest), macOS (latest)
    \item \textbf{Python Versions}: 3.10, 3.11, 3.12
    \item \textbf{Test Matrix}: 9 combinations (3 OS $\times$ 3 Python versions)
\end{itemize}

\paragraph{Quick Validation:}
A fast smoke test workflow (\filepath{.github/workflows/quick-test.yml}) runs on every commit to provide rapid feedback during development, using Ubuntu with Python 3.11 for minimal latency.

\paragraph{Code Coverage:}
Test coverage is tracked using \texttt{pytest-cov}, with HTML reports generated for detailed inspection of tested code paths.

All tests must pass before pull requests are merged, ensuring the pipeline remains stable across platforms and Python versions. The current test suite achieves complete coverage of critical configuration and import paths, with plans to expand coverage to include data processing functions.

\subsection{Reproducibility and Version Control}

To ensure full reproducibility, the repository provides:

\paragraph{Dependency Management:}
\begin{itemize}
    \item \filepath{requirements.txt}: Flexible dependency specifications for active development
    \item \filepath{requirements-frozen.txt}: Exact pinned versions for bit-for-bit reproduction
    \item \filepath{environment.yml}: Conda environment specification with channel information
\end{itemize}

\paragraph{Citation Metadata:}
\begin{itemize}
    \item \filepath{CITATION.cff}: Machine-readable citation file following Citation File Format standard
    \item Includes DOI placeholders for both code and paper
    \item Compatible with GitHub's built-in citation feature
    \item Provides BibTeX export automatically
\end{itemize}

\paragraph{Version Control:}
\begin{itemize}
    \item Semantic versioning (v1.0.0 format)
    \item Tagged releases on GitHub
    \item Comprehensive commit history with descriptive messages
    \item Protected main branch requiring review before merging
\end{itemize}

\subsection{Future Extensions}

The modular architecture facilitates several planned enhancements:

\paragraph{Automated Updates:} Implementation of scheduled execution (e.g., monthly cron jobs) to automatically incorporate new Weekly Report publications.

\paragraph{Alternative Data Sources:} Extension to INEI's direct publications or other statistical agencies with minimal code modifications due to configuration-driven design.

\paragraph{Additional Frequencies:} Support for weekly or daily economic indicators using the same pipeline architecture.

\paragraph{Statistical Validation:} Integration of automated rationality tests and revision analysis directly in the pipeline, beyond the current smoke tests.

\paragraph{API Development:} Creation of RESTful API for programmatic access to RTD without local installation, enabling integration with other systems.

\paragraph{Real-Time Updates:} Deployment of the dashboard as a public web service with automatic monthly updates as new Weekly Reports are published.

The codebase design explicitly supports these extensions through:
\begin{itemize}
    \item Modular separation of concerns (scrapers, processors, cleaners, transformers)
    \item Configuration-driven execution (new sources via YAML, not code changes)
    \item Comprehensive type hints (enables IDE support and type checking)
    \item Extensive documentation (docstrings, README, technical supplement)
    \item Version control (Git, semantic versioning)
\end{itemize}

\subsection{Code Availability and Support}

\paragraph{Public Repository:} \url{https://github.com/JasonCruz18/peruvian_gdp_revisions}

\paragraph{License:} MIT License (permits commercial and academic use, modification, and redistribution)

\paragraph{Documentation:}
\begin{itemize}
    \item Repository README with installation and usage instructions
    \item Comprehensive \filepath{docs/} directory with guides for installation, usage, architecture, and contribution
    \item Six tutorial Jupyter notebooks covering all pipeline stages
    \item FAQ document addressing common questions and troubleshooting
    \item Inline docstrings for all functions and classes
    \item Complete type hints throughout codebase
    \item Configuration file comments and examples
    \item This technical supplement
\end{itemize}

\paragraph{Support Channels:}
\begin{itemize}
    \item GitHub Issues: Bug reports and feature requests
    \item GitHub Discussions: General questions and usage help
    \item Email: \href{mailto:jj.cruza@up.edu.pe}{jj.cruza@up.edu.pe} for research collaboration
\end{itemize}

\paragraph{Citation:} Researchers using this code or dataset should cite both the paper and the repository:

\begin{lstlisting}[language=tex, caption={Suggested citation}]
@article{CruzWinkelriedTorres2025gdp,
  title={Nowcasting GDP using data revisions in an emerging economy},
  author={Cruz, Jason and Winkelried, Diego and Torres, Javier},
  year={2025},
  institution={Universidad del Pacifico - CIUP}
}

@software{CruzWinkelriedTorres2025gdp_code,
  title={Peru GDP Real-Time Dataset Construction Pipeline},
  author={Cruz, Jason},
  year={2025},
  url={https://github.com/JasonCruz18/peruvian_gdp_revisions},
  version={1.0.0}
}
\end{lstlisting}

The combination of comprehensive documentation, modular design, and public availability ensures that this computational work meets the highest standards of reproducibility. All materials necessary for complete replication are freely available, and the codebase is designed to facilitate both exact replication and extensions for future research.

\end{document}
